\documentclass[11pt,a4paper]{article}

\usepackage[a4paper,margin=2.2cm]{geometry}
\usepackage[T1]{fontenc}
\usepackage[utf8]{inputenc}
\usepackage{lmodern}
\usepackage{amsmath,amssymb}
\usepackage{array}
\usepackage{booktabs}
\usepackage{enumitem}
\usepackage{longtable}
\usepackage{graphicx}
\usepackage{pdflscape}
\usepackage{titlesec}
\usepackage{url}

\setlength{\parindent}{0pt}
\setlength{\parskip}{0.6em}

\setlist[itemize]{leftmargin=1.5em,itemsep=0.2em,topsep=0.3em,parsep=0pt,partopsep=0pt}
\setlist[enumerate]{leftmargin=1.8em,itemsep=0.2em,topsep=0.3em,parsep=0pt,partopsep=0pt}

\titlespacing*{\subsection}{0pt}{0.9em}{0.4em}

\newcommand{\reportsection}[1]{%
  \refstepcounter{section}%
  \addcontentsline{toc}{section}{\protect\numberline{\thesection}#1}%
  \vspace{1.2em}%
  \noindent\rule{\textwidth}{0.4pt}\par\vspace{0.8ex}%
  \noindent{\Large\bfseries \thesection\quad #1}\par\vspace{0.5em}%
}

\begin{document}

\begin{titlepage}
\thispagestyle{empty}
\begin{center}
\fontfamily{ptm}\selectfont
\fontsize{12}{24}\selectfont

King's College London

\vspace*{0.22\textheight}

Comparative Evaluation of Optimisation Approaches for Hospital Resource Scheduling

\vspace*{0.16\textheight}

Maksym Byelko \& Kexin Zhong\\
EPSRC Summer Internship\\
Steffen Zschaler\\
August 2026
\end{center}
\end{titlepage}

\tableofcontents
\clearpage

\reportsection{Aim}

This experiment evaluates an MDEO-Cloud implementation against a Python-based optimiser for the International Hospital Timetabling Competition (IHTC) scheduling problem~\cite{ceschia2025ihtc}. Both solve the same integrated hospital-planning problem. Patients are assigned admission days, rooms, and operating theatres, after which nurses are allocated to the room-day-shift demand created by those assignments.

The evaluation examines feasibility, schedule quality, and the computational effort needed to obtain a schedule. Each approach uses the same IHTC benchmark instances and input data. MDEO-Cloud executions use a fixed limit of 300 generations, regardless of whether the final candidate set has converged. Runtime results must also account for execution conditions, including the different thread allocation used for the MDEO-Cloud Test~10 run.

The study records patient-admission outcomes, resource use, nursing coverage, constraint violations, and search-performance measures. It compares the results obtained on the selected benchmarks and considers the practical implications of using model-driven optimisation in MDEO-Cloud instead of a conventional Python implementation for hospital scheduling.

\reportsection{Controlled Variables}

Both implementations use the same IHTC benchmark instance set and associated input data. This holds the scheduling problem constant while the implementation approach varies. Runtime, feasibility, and solution quality are recorded as experimental outcomes.

\begin{itemize}
    \item \textbf{Patients and occupants:} patient priority, release and due dates, length of stay, surgery duration, assigned surgeon, room compatibility, and existing occupants.
    \item \textbf{Nursing data:} nurse skill levels, working-shift availability, maximum workload, and daily nursing demand for each patient.
    \item \textbf{Surgical resources:} surgeon availability and operating-theatre capacity for each day.
    \item \textbf{Room resources:} room capacity, compatibility requirements, and initial occupancy.
    \item \textbf{Planning structure:} the decision horizon, available rooms, shift definitions, and scheduling rules specified by the IHTC instance.
\end{itemize}

\reportsection{Problem Definition}

The IHTC problem links patient scheduling with nurse assignment. Patient scheduling determines whether a patient is admitted and, where admission occurs, assigns an admission day, room, and operating theatre. The selected date must lie within the patient's release and due dates. It must also respect the availability of the assigned surgeon and operating theatre.

Room constraints apply throughout a patient's stay. The assigned room must have capacity on every occupied day, and patients in different age groups may not share a room on the same day. Existing occupants use room capacity and create nursing demand, although they are not scheduling decisions.

Active room-day-shift slots are derived from admitted patients and their daily demand records. A nurse may be assigned to a slot when they work that day and shift and meet its skill requirement. The resulting allocation is assessed using coverage, uncovered workload, workload excess, and workload balance~\cite{ceschia2025ihtc}.

The metamodel captures the relevant decisions and resource states. It includes patients, admissions, rooms and room-day availability, surgeons and their availability, operating theatres and their availability, nurses and working shifts, patient-day demands, hospitalisation shifts, and room-shift assignments. Optimisation transformations revise candidate schedules, while objective functions evaluate them.

\reportsection{Dataset Sizes}

The evaluation uses ten IHTC 2024 benchmark instances. They differ in the number of patients awaiting admission, existing occupants, and available hospital resources. Existing occupants are reported separately because they occupy rooms and generate nursing demand, although they are not admission decisions.

\begin{table}[htbp]
\centering
\small
\begin{tabular}{lrrrrrrr}
\toprule
\textbf{Test} & \textbf{Days} & \textbf{Patients} & \textbf{Occupants} & \textbf{Nurses} & \textbf{Rooms} & \textbf{Surgeons} & \textbf{Theatres} \\
\midrule
IHTC 2024 Test 01 & 21 & 42 & 7 & 13 & 5 & 1 & 2 \\
IHTC 2024 Test 02 & 14 & 37 & 5 & 17 & 6 & 2 & 3 \\
IHTC 2024 Test 03 & 14 & 45 & 10 & 14 & 6 & 1 & 2 \\
IHTC 2024 Test 04 & 14 & 54 & 7 & 19 & 8 & 2 & 3 \\
IHTC 2024 Test 05 & 14 & 62 & 9 & 15 & 6 & 1 & 2 \\
IHTC 2024 Test 06 & 14 & 111 & 9 & 20 & 9 & 3 & 3 \\
IHTC 2024 Test 07 & 21 & 113 & 13 & 22 & 9 & 2 & 3 \\
IHTC 2024 Test 08 & 21 & 173 & 15 & 21 & 9 & 2 & 3 \\
IHTC 2024 Test 09 & 21 & 146 & 26 & 29 & 14 & 2 & 2 \\
IHTC 2024 Test 10 & 21 & 525 & 50 & 83 & 46 & 10 & 11 \\
\bottomrule
\end{tabular}
\caption{Resource composition of the IHTC 2024 benchmark instances~\cite{ceschia2025ihtc}. Values were derived from the supplied JSON benchmark files.}
\end{table}

Test labels do not rank the instances by size. For example, Test 02 contains fewer patients than Test 01. The benchmark set therefore examines scalability across irregular instance sizes rather than fixed artificial categories.

\reportsection{Optimisation Approaches}

Both implementations should be evaluated using the same instance set. MDEO-Cloud runs terminate after 300 generations. Runtime comparisons require equivalent execution resources.

\subsection{MDEO-Cloud Solution}

The MDEO-Cloud implementation encodes each IHTC instance as a typed model. Patients, admissions, resources, availability records, nursing-demand records, and nurse assignments are represented as explicit model elements. During conversion, every mandatory patient is initialised as scheduled and linked to an admission; optional patients remain unscheduled. Each candidate schedule is stored as a complete model instance and processed within the MDEO-Cloud platform~\cite{krieger2026mdeo}.

NSGA-II~\cite{deb2002nsga2} explores candidate schedules through model transformations. These transformations add, remove, or reassign patient admissions, rooms, admission days, and operating theatres. They also create, remove, and move nurse assignments between room-day-shift slots. Guards reject infeasible changes before the transformations are applied. A phase state transfers the search from patient scheduling to nurse assignment after all mandatory patients have been admitted.

Objective functions calculate resource use and nursing demand from the current admission assignments and patient-day demand records. The MDEO-Cloud implementation optimises the following objectives:

\begin{itemize}
    \item maximise admitted patients.
    \item minimise the unscheduled-patient penalty, which applies a higher penalty to mandatory patients, and minimise admission delay.
    \item minimise room-capacity violation and missing mandatory-patient admissions as redundant safeguard objectives. These do not replace hard enforcement: room capacity is guarded by transformations, while the configured phase-consistency constraint requires every mandatory patient to be admitted in a final candidate.
    \item minimise aggregate surgeon and operating-theatre operating-time excess.
    \item minimise uncovered room-shifts and nursing workload.
    \item minimise excess and imbalanced nurse workload.
    \item minimise inactive nurse assignments.
    \item minimise room-occupancy standard deviation.
\end{itemize}

\subsection{Python-Based Solution}
The Python-based implementation formulates the IHTC scheduling problem as a Mixed Integer Linear Programming (MILP) model built upon the PuLP optimisation library. This exact optimisation framework differs fundamentally from the heuristic NSGA-II evolutionary algorithm adopted in MDEO-Cloud. All hard constraints laid down by the IHTC competition regulations are explicitly encoded as linear equalities and inequalities to guarantee the feasibility of generated schedules.

The objective function adopts a weighted single-objective minimisation structure which aggregates eight soft penalty terms aligned with the competition evaluation metrics, corresponding to soft constraints S1 to S8: unscheduled patient penalty, admission delay penalty, open operating theatre penalty, age-mixing penalty, skill mismatch penalty, excess nurse workload penalty, continuity-of-care penalty, and surgeon transfer penalty. Weight parameters for each cost component follow the official benchmark configuration to ensure a fair comparison across the two solution approaches.

The model comprises two sequential decision layers: patient admission and room allocation, followed by nurse assignment for room-day-shift slots. Binary decision variables are defined to represent patient admission status, room occupancy, operating theatre activation and nurse-room shift allocation. The Big-M linearisation technique is utilised to convert non-linear logical rules into standard linear constraints. PuLP invokes the open-source CBC solver to find optimal or near-optimal feasible solutions under pre-defined computational time limits.

In contrast to the iterative evolutionary search of MDEO-Cloud, the MILP approach can deliver mathematically provable optimal solutions given sufficient computational resources. Nevertheless, the quantity of variables and constraints expands rapidly as the instance scale grows, imposing substantial memory and runtime overhead for larger benchmark test cases.

\reportsection{Constraints}

\subsection{MDEO-Cloud Solution}

The MDEO-Cloud solution uses one configuration-level constraint and a set of transformation guards. The configuration-level constraint enforces completion of the nurse-assignment phase and admission of every mandatory patient. Guards enforce the remaining scheduling and nurse-assignment rules before a transformation is applied. They check that an individual admission fits the selected surgeon and theatre availability records. Aggregate surgeon and theatre operating-time excess across all patients on a day is evaluated by the objective functions.

\begin{itemize}
    \item \textbf{Phase consistency:} the \texttt{nurseAssignmentPhaseConsistency} constraint returns a non-zero violation for a final candidate in the \texttt{PATIENTS} phase. In the \texttt{NURSES} phase, it returns the number of mandatory patients without an admission. A valid final candidate must therefore reach \texttt{NURSES} and admit every mandatory patient. The \texttt{mandatoryPatientsAdmitted()} function is a helper and is not configured as a constraint.

    \item \textbf{Admission scope and uniqueness:} mandatory admissions are created during model conversion rather than by optimisation transformations. Transformations T1--T3 add only unscheduled, non-mandatory patients and prevent duplicate admissions for the selected patient. T4 removes optional patients only, T5 and T7 may remove any scheduled patient, and T6 replaces optional patients. A final model missing a mandatory admission remains infeasible through \texttt{nurseAssignmentPhaseConsistency}.

    \item \textbf{Patient admission window:} a new admission day must lie between the patient's release date and due date, and the complete patient stay must remain within the decision horizon.

    \item \textbf{Surgeon and theatre feasibility:} admission, day-move, and theatre-move transformations link each patient to their assigned surgeon and require the individual surgery duration to fit the selected surgeon availability record. They also require a theatre availability record with positive capacity on the chosen day. Aggregate daily capacity excess is evaluated as an objective.

    \item \textbf{Room compatibility:} a patient cannot be assigned to a room listed in their incompatible rooms.

    \item \textbf{Full-stay room capacity:} every room-day within the patient's stay must have spare capacity before an admission or room reassignment is applied.

    \item \textbf{Room age-group compatibility:} each occupied room-day must either be empty or already contain patients from the same age group. Patients from different age groups cannot be mixed in the same room on the same day.

    \item \textbf{Room-state consistency:} admission, removal, and reassignment transformations update \texttt{occupiedBeds} for every room-day in the stay and maintain \texttt{AgeGroup.EMPTY} when a room-day becomes unoccupied.

    \item \textbf{Transformation phase control:} transformations T4--T13 operate during the \texttt{PATIENTS} phase, and T14--T16 require \texttt{NURSES}. T1--T3 may add optional patients in either phase. The one-way transition to \texttt{NURSES} is available only after every mandatory patient has been scheduled.

    \item \textbf{Nurse availability:} a nurse can only be assigned to a room-day-shift when they have a matching \texttt{NurseWorkingShift} on the same day and shift.

    \item \textbf{Nurse skill compatibility:} a nurse assignment is rejected if any patient demand contributing to the target room-day-shift requires a higher skill level than that nurse possesses.

    \item \textbf{Active nurse assignments:} nurses can only be assigned to room-day-shift slots for which a current admission and patient-day demand generate nursing demand. Duplicate assignments to the same slot are prevented.

    \item \textbf{Objective-based penalties:} the unscheduled-patient penalty, room-capacity violation, missing mandatory-patient admissions, aggregate surgeon and theatre operating-time excess, nurse workload excess, uncovered room-shifts, uncovered workload, workload imbalance, inactive nurse assignments, admission delay, and room-occupancy standard deviation are minimised as objective values. The room-capacity and missing-mandatory-admission objectives are redundant safeguards alongside their hard enforcement.
\end{itemize}

\subsection{Python-Based Solution}
The Python-based MILP model strictly separates hard constraints (which cannot be violated under any feasible solution) and soft penalty constraints (permitted to be violated but incur additive cost within the objective function). All hard rules defined by the IHTC competition are encoded as linear equality and inequality constraints enforced by the solver. No constraint violation is allowed in any feasible schedule.

\begin{itemize}
    \item \textbf{Patient admission window}: The admission day of each scheduled patient must lie within the interval between the patient’s release date and due date. The complete length-of-stay period must not exceed the planning horizon.
    \item \textbf{Surgeon and theatre feasibility}: A patient may only be admitted on a day where their assigned surgeon has available operating capacity and the selected operating theatre possesses residual capacity on that day.
    \item \textbf{Room compatibility}: A patient cannot be allocated to any room marked within their incompatible room list.
    \item \textbf{Full-stay room capacity}: For every day covered by the patient’s stay, the assigned room must maintain spare capacity to accommodate the patient.
    \item \textbf{Room age-group compatibility}: Patients from distinct age groups cannot occupy the same room on the same day. This rule is implemented as a hard constraint.
    \item \textbf{Nurse availability}: A nurse can only be assigned to a room-day-shift slot if they hold a valid \textit{NurseWorkingShift} record for that exact day and shift.
    \item \textbf{Nurse skill compatibility}: The skill level of an assigned nurse must satisfy or exceed the maximum skill demand generated by patients within the corresponding room-day-shift.
    \item \textbf{Active nurse assignments}: Nurse allocation is only enabled for room-day-shift entries with active patient admissions and positive nursing demand. Multiple assignments of the same nurse to one room-day-shift are prohibited.
\end{itemize}

Resource excess conditions, age-group mixing, uncovered nursing workload, nurse workload overload, care continuity degradation and other performance metrics are not implemented as hard constraints. These rules correspond to soft constraints S1–S8, and all associated penalties are aggregated into the weighted objective function for minimisation.

\reportsection{Evaluation Metrics}

The experiment evaluates solutions in two stages: feasibility first, followed by solution quality and search performance. For MDEO-Cloud, feasibility is an implementation-derived classification rather than the result of a separate post-run validator. A final MDEO-Cloud model is feasible when \texttt{nurseAssignmentPhaseConsistency} returns zero, which requires the model to be in the \texttt{NURSES} phase with every mandatory patient admitted, and when its state has been produced through transformations whose guards preserve the specified hard scheduling invariants. If no MDEO-Cloud model meets this classification after the fixed 300-generation limit, the run is classified as infeasible and excluded from feasible-solution quality summaries. Its diagnostic and objective values are reported instead.

\subsection{MDEO-Cloud Solution}

\subsubsection{Feasibility Metrics}

\begin{itemize}
    \item Final feasibility status: whether the reported final candidate satisfies the MDEO-Cloud feasibility classification: zero \texttt{nurseAssignmentPhaseConsistency} and a state reached through the guarded transformations.
    \item Reported hard diagnostic: room-capacity status derived from the final model state. Admission windows, incompatible rooms, age-group compatibility, and nurse-assignment rules are enforced by transformation guards during search; they are not independently audited as post-run violation counts. Theatre and surgeon operating-time excess, together with uncovered workload, are tabulated as objective values. Nurse workload excess is retained only in the raw Pareto-front outcomes.
    \item Mandatory-patient completion rate: proportion of mandatory patients that receive an admission.
\end{itemize}

\subsubsection{Patient Scheduling Metrics}

\begin{itemize}
    \item Number and proportion of optional patients admitted.
    \item Number of unscheduled optional patients.
    \item Total and mean admission delay, measured from each patient's release date to their admission day.
    \item Number of open operating theatres per day.
    \item Number of theatres used by each surgeon per day.
    \item Room age-group compatibility: enforced by transformation guards during search and not reported as a post-run violation count.
\end{itemize}

\subsubsection{Nurse Assignment Metrics}

The stored Pareto-front outcomes retain objective values for uncovered room-shifts, uncovered workload, nurse workload excess, and nurse workload imbalance. Of these, only uncovered workload is presented in the final results table; the remaining values are available in the raw MDEO-Cloud run outputs but are not tabulated in this report.

\begin{itemize}
    \item \textbf{Reported metric --- uncovered workload:} total workload associated with active room-shift slots without nurse cover.
    \item \textbf{Raw outcome metrics --- uncovered room-shifts, nurse workload excess, and nurse workload imbalance:} retained in the Pareto-front objective values but not tabulated in the final results.
    \item \textbf{Nurse skill compatibility:} enforced by transformation guards. Although \texttt{totalSkillLevelViolation()} is defined and imported, it is not configured as a goal objective or formal constraint and is not reported as a final numeric metric.
\end{itemize}

\subsubsection{Reported Resource Measures}

\begin{itemize}
    \item \textbf{Reported hard diagnostic --- room capacity:} the final room-capacity status shown in Table~\ref{tab:feasibility-results}.
    \item \textbf{Reported objective measures --- theatre and surgeon operating-time excess:} aggregate excess values shown in Table~\ref{tab:feasibility-results}.
\end{itemize}

Room-bed, theatre, and surgeon utilisation ratios were not exported or reported for the completed MDEO-Cloud runs.

\subsubsection{Search Performance Metrics}

For each MDEO-Cloud execution, the reported search measures are:

\begin{itemize}
    \item wall-clock runtime, completed generations, and candidate evaluations;
    \item executed and skipped transformations, and the number of Pareto solutions;
    \item final feasibility status and the final model's objective and diagnostic values;
\end{itemize}

\subsubsection{Reporting}

Results are reported as a set of feasibility, patient-scheduling, nurse-assignment, resource-related, and search-performance measures. The evaluation does not combine these measures into a single score. Some constraints are represented as hard conditions in one implementation and as soft penalties in the other, so implementation-specific objective values and penalty totals are interpreted in their own context. MDEO-Cloud raw outcomes include missing mandatory admissions, inactive nurse assignments, room-occupancy standard deviation, and transformation skip rate.

\subsection{Python-Based Solution}
Evaluation of the MILP implementation follows the same two-stage framework: feasibility verification followed by solution quality assessment. A candidate schedule is only regarded as valid if all hard constraints are fully satisfied without violation. All feasible solutions are then evaluated using the unified official IHTC weighted cost function, which aggregates the eight soft penalty components S1–S8 as the primary comparison metric.

The evaluation metrics are categorised consistently with the MDEO-Cloud benchmark, covering feasibility statistics, patient scheduling outcomes, nurse assignment performance and resource utilisation. Metrics specific to evolutionary search, such as generation count and transformation attempts, are omitted, as they are not applicable to the exact MILP solving paradigm.

Solver performance is recorded via wall-clock runtime and final solver status (optimal, feasible or time-limited). Detailed per-instance soft penalty breakdowns are recorded to support cross-method comparison against results obtained from MDEO-Cloud.

\reportsection{System Hardware}

The MDEO-Cloud experiments were executed on a system with an AMD Ryzen 9 9950X3D processor and 96\,GB of DDR5 RAM. Tests~01--09 ran on a development virtual machine hosted on this hardware, with a smaller allocation of CPU threads; the precise allocation was not retained. Test~10 used three MDEO execution nodes with eight threads each (24 threads in total) because the earlier allocation would have made the execution impractically long. Its runtime is therefore an annotated large-instance observation and is not directly comparable with Tests~01--09 as a like-for-like scalability result.

For Test~10, the MDEO-Cloud script timeout and model-transformation timeout were each set to 60 seconds as safeguards. Neither threshold was reached, so these settings did not alter the run termination or its reported result.

\reportsection{Experiment Procedure}

Each benchmark instance is executed once per framework. For MDEO-Cloud, Tests~01--09 used the development virtual machine allocation, while Test~10 used additional CPU threads. The configured Test~10 safeguards were not reached.

\subsection{MDEO-Cloud Solution}

For the completed MDEO-Cloud executions, the workflow was:

\begin{itemize}
    \item \textbf{Instance Collection:} the IHTC test instances were collected for the experiment.
    \item \textbf{Model Conversion:} each test instance was converted into a valid MDEO-Cloud model representation.
    \item \textbf{Model Validation:} the converted models were checked for required entities, links, and availability objects.
    \item \textbf{Configuration:} solver hyperparameters were set, including population size, number of evolutions, and mutation settings.
    \item \textbf{Termination:} MDEO-Cloud used no whole-run wall-clock limit. Every run terminated after 300 generations, regardless of whether its candidate solutions had converged. Wall-clock duration was recorded as an experimental result. Test~10 also had 60-second script and model-transformation safeguards, but neither threshold was reached and they did not terminate the run.
    \item \textbf{Optimisation Execution:} the framework was run with the implemented metamodel, transformations, constraints, and objective functions.
    \item \textbf{Schedule Generation:} final model solutions, including admissions and room-shift nurse assignments, were obtained.
    \item \textbf{Result Export:} final schedules were exported or converted into a readable format for analysis.
    \item \textbf{Data Collection:} objective values, violations, runtime behaviour, and final schedule statistics were recorded for comparison.
\end{itemize}

\subsection{Python-Based Solution}

\textit{To be completed by the author of the Python-based solution.}\par

\reportsection{Quantitative Analysis}

\subsection{MDEO-Cloud Solution}

Results are reported separately for each benchmark instance. Feasibility, patient admissions, nursing coverage and workload, resource utilisation, runtime, and search-efficiency measures provide the quantitative evidence. The evaluation does not reduce these measures to a single score, and implementation-specific objective values remain in their own context. Results are not pooled across instances because the instances differ substantially in scale.

For each instance, the following values are recorded where available:

\begin{itemize}
    \item final feasibility status and final objective and diagnostic values;
    \item wall-clock runtime, candidate evaluations, and completed generations;
    \item patient-admission, nursing-coverage, resource-utilisation, and constraint-diagnostic values.
\end{itemize}

These measures support the comparison of optimisation quality, runtime, and scalability between the two frameworks.

\subsection{Python-Based Solution}

\textit{To be completed by the author of the Python-based solution.}\par

\reportsection{Results}

\subsection{Per-Instance Comparison}

Table~\ref{tab:instance-results} reports per-instance feasibility and runtime. The associated feasibility, coverage, resource, and runtime measures are reported separately.

\begin{table}[htbp]
\centering
\scriptsize
\resizebox{\textwidth}{!}{%
\begin{tabular}{lrrrrr}
\toprule
\textbf{Instance} & \textbf{Patients} & \textbf{MDEO Feasible} & \textbf{MDEO Time (s)} & \textbf{Python Feasible} & \textbf{Python Time (s)} \\
\midrule
IHTC 2024 Test 01 & 42  & Yes & 521  & -- & -- \\
IHTC 2024 Test 02 & 37  & Yes & 363  & -- & -- \\
IHTC 2024 Test 03 & 45  & Yes & 363  & -- & -- \\
IHTC 2024 Test 04 & 54  & Yes & 647  & -- & -- \\
IHTC 2024 Test 05 & 62  & Yes & 568  & -- & -- \\
IHTC 2024 Test 06 & 111 & Yes & 1,454 & -- & -- \\
IHTC 2024 Test 07 & 113 & Yes & 2,240 & -- & -- \\
IHTC 2024 Test 08 & 173 & Yes & 3,282 & -- & -- \\
IHTC 2024 Test 09 & 146 & Yes & 3,728 & -- & -- \\
IHTC 2024 Test 10$^\ast$ & 525 & Yes & 32,676 & -- & -- \\
\bottomrule
\end{tabular}%
}
\caption{Per-instance feasibility and runtime results. $^\ast$Test~10 used additional CPU threads; its runtime is an annotated large-instance observation and is not directly comparable with Tests~01--09.}
\label{tab:instance-results}
\end{table}

\subsection{Feasibility and Constraint Results}

Table~\ref{tab:feasibility-results} distinguishes the implementation-derived feasibility classification and reported room-capacity diagnostic from objective measures. Theatre operating-time excess, surgeon operating-time excess, and uncovered nursing workload are objective values. A feasible candidate may therefore have non-zero values for these objective measures.

\begin{landscape}
{\scriptsize
\setlength{\tabcolsep}{3pt}
\begin{longtable}{llrrrrrrr}
\caption{Feasibility, hard room-capacity diagnostics, and reported objective measures by framework.}\label{tab:feasibility-results}\\
\toprule
\textbf{Instance} & \textbf{Framework} & \textbf{Feasible} & \textbf{Mandatory} & \textbf{All Patients} & \textbf{Room Capacity} & \textbf{Theatre Excess} & \textbf{Surgeon Excess} & \textbf{Uncovered} \\
& & & \textbf{Admitted} & \textbf{Admitted} & \textbf{(Hard)} & \textbf{(Objective)} & \textbf{(Objective)} & \textbf{Workload (Objective)} \\
\midrule
\endfirsthead
\toprule
\textbf{Instance} & \textbf{Framework} & \textbf{Feasible} & \textbf{Mandatory} & \textbf{All Patients} & \textbf{Room Capacity} & \textbf{Theatre Excess} & \textbf{Surgeon Excess} & \textbf{Uncovered} \\
& & & \textbf{Admitted} & \textbf{Admitted} & \textbf{(Hard)} & \textbf{(Objective)} & \textbf{(Objective)} & \textbf{Workload (Objective)} \\
\midrule
\endhead
IHTC 2024 Test 01 & MDEO-Cloud & Yes & 11/11 & 31/42 & 0 & 0 & 0 & 959 \\
IHTC 2024 Test 01 & Python & -- & -- & -- & -- & -- & -- & -- \\
IHTC 2024 Test 02 & MDEO-Cloud & Yes & 12/12 & 32/37 & 0 & 0 & 0 & 1,034 \\
IHTC 2024 Test 02 & Python & -- & -- & -- & -- & -- & -- & -- \\
IHTC 2024 Test 03 & MDEO-Cloud & Yes & 12/12 & 18/45 & 0 & 0 & 0 & 352 \\
IHTC 2024 Test 03 & Python & -- & -- & -- & -- & -- & -- & -- \\
IHTC 2024 Test 04 & MDEO-Cloud & Yes & 32/32 & 50/54 & 0 & 0 & 0 & 1,418 \\
IHTC 2024 Test 04 & Python & -- & -- & -- & -- & -- & -- & -- \\
IHTC 2024 Test 05 & MDEO-Cloud & Yes & 19/19 & 26/62 & 0 & 0 & 0 & 603 \\
IHTC 2024 Test 05 & Python & -- & -- & -- & -- & -- & -- & -- \\
IHTC 2024 Test 06 & MDEO-Cloud & Yes & 50/50 & 68/111 & 0 & 0 & 0 & 1,772 \\
IHTC 2024 Test 06 & Python & -- & -- & -- & -- & -- & -- & -- \\
IHTC 2024 Test 07 & MDEO-Cloud & Yes & 55/55 & 67/113 & 0 & 0 & 0 & 2,014 \\
IHTC 2024 Test 07 & Python & -- & -- & -- & -- & -- & -- & -- \\
IHTC 2024 Test 08 & MDEO-Cloud & Yes & 39/39 & 77/173 & 0 & 0 & 0 & 2,017 \\
IHTC 2024 Test 08 & Python & -- & -- & -- & -- & -- & -- & -- \\
IHTC 2024 Test 09 & MDEO-Cloud & Yes & 88/88 & 98/146 & 0 & 0 & 0 & 2,960 \\
IHTC 2024 Test 09 & Python & -- & -- & -- & -- & -- & -- & -- \\
IHTC 2024 Test 10 & MDEO-Cloud & Yes & 349/349 & 349/525 & 0 & 0 & 0 & 9,289 \\
IHTC 2024 Test 10 & Python & -- & -- & -- & -- & -- & -- & -- \\
\bottomrule
\end{longtable}
}
\end{landscape}

\subsection{Performance and Scalability}

\subsubsection{MDEO-Cloud Solution}

Tests~01--10 each terminated after the fixed limit of 300 generations. Each run recorded 59,900 candidate evaluations, which is an observed search measure rather than an additional termination limit. Tests~01--09 used the same development virtual machine allocation and provide the internally comparable runtime trend: the smaller cases with 37--62 patients completed in 363--647~s, while cases with 111--173 patients took 1,454--3,728~s. Test~10, which has 525 patients, used additional CPU threads and recorded 32,676~s (approximately 9~hours and 5~minutes). Its 60-second script and model-transformation safeguards were never reached. Test~10 is therefore retained as an annotated large-instance observation, rather than a directly comparable measurement of the Tests~01--09 scalability trend.

Runtime was not determined by patient count alone: Test~09, with 146 patients, took longer than Test~08, with 173 patients. This difference is consistent with the variation in executed and skipped transformations shown in Table~\ref{tab:mdeo-diagnostics}, but a single execution per instance cannot isolate the effect of any individual instance characteristic. The same caution applies more strongly to Test~10: it has substantially more candidate patients but different execution resources, so its recorded runtime cannot be attributed to patient count or instance structure alone. The unused safeguards did not terminate or otherwise alter the run.

For Test~10, the 300-generation budget did not show clear convergence. At the final generation, the search was still applying transformations and retained substantial trade-offs between patient admissions and violation-related measures. The reported 349-admission schedule should therefore not be read as evidence that the search had reached its best attainable schedule.

No single score is calculated from the MDEO schedules reported in the results. No mean, standard deviation, or statistical significance claim is reported because each instance was executed once.

\subsubsection{Python-Based Solution}

\textit{To be completed by the author of the Python-based solution. Record corresponding runtime and IHTC-score results once available.}

\subsection{MDEO-Cloud Diagnostics}

Table~\ref{tab:mdeo-diagnostics} records MDEO-specific search behaviour. These values explain MDEO performance but are not used directly to compare the Python-based implementation.

\begin{table}[htbp]
\centering
\scriptsize
\resizebox{\textwidth}{!}{%
\begin{tabular}{lrrrrr}
\toprule
\textbf{Instance} & \textbf{Generations} & \textbf{Candidate Evaluations} & \textbf{Executed Transformations} & \textbf{Skipped Transformations} & \textbf{Pareto Solutions} \\
\midrule
IHTC 2024 Test 01 & 300 & 59,900 & 113,202 & 54,708 & 100 \\
IHTC 2024 Test 02 & 300 & 59,900 & 123,565 & 79,345 & 100 \\
IHTC 2024 Test 03 & 300 & 59,900 & 110,520 & 47,546 & 100 \\
IHTC 2024 Test 04 & 300 & 59,900 & 120,264 & 71,268 & 100 \\
IHTC 2024 Test 05 & 300 & 59,900 & 118,583 & 56,472 & 100 \\
IHTC 2024 Test 06 & 300 & 59,900 & 109,820 & 50,207 & 100 \\
IHTC 2024 Test 07 & 300 & 59,900 & 138,207 & 131,048 & 100 \\
IHTC 2024 Test 08 & 300 & 59,900 & 120,323 & 54,716 & 100 \\
IHTC 2024 Test 09 & 300 & 59,900 & 118,060 & 63,923 & 100 \\
IHTC 2024 Test 10 & 300 & 59,900 & 140,014 & 94,769 & 100 \\
\bottomrule
\end{tabular}%
}
\caption{MDEO-Cloud search diagnostics.}
\label{tab:mdeo-diagnostics}
\end{table}

\subsection{Python-Based Diagnostics}

\textit{To be completed by the author of the Python-based solution. Record only framework-specific diagnostics that are available and useful for explaining performance.}

\reportsection{Interpretation}

\subsection{Cross-Implementation Interpretation}

Interpretation should be limited to the evidence in Tables~\ref{tab:instance-results} and~\ref{tab:feasibility-results}. It should describe how feasibility, admissions, nursing coverage and workload, resource use, and runtime differed across completed instances, together with the constraints or search behaviours that limited each approach. Test~10 runtime should be interpreted only in the context of its additional CPU threads, and should not be compared directly with another framework unless execution conditions are equivalent. Full schedules, raw logs, Pareto fronts, and Excel exports should remain in the project repository or appendix.

\subsection{Comparison with Existing GIPS-Based Results}

The literature also provides a useful external reference point through the recent GIPS-based MILP implementation reported for the IHTC 2024 problem~\cite{kratz2025gips}. That study uses the same benchmark family but not the same experimental setup as the present report: it evaluates all 30 public IHTC instances, whereas the current MDEO-Cloud evaluation reports Tests~01--10 only. The published GIPS approach also relies on Gurobi within a MILP pipeline, rather than NSGA-II search in MDEO-Cloud.

For that reason, the GIPS results should be interpreted as contextual benchmark evidence rather than as a third arm of the controlled experiment. The GIPS paper reports feasible solutions for all 30 public instances within the original 10-minute competition limit and proves optimality for two public instances under a relaxed runtime budget. The present MDEO-Cloud study reports feasible schedules for all 10 tested instances. These outcomes describe different experimental settings and do not establish a direct ranking of runtime or benchmark coverage. Within the MDEO-Cloud study, Tests~01--09 provide the internally comparable runtime observations; Test~10 is not used for contextual runtime comparison because it had a different CPU-thread allocation.

\reportsection{Qualitative Analysis}

In addition to quantitative results, the experiment will include a qualitative comparison of the two frameworks based on implementation and debugging experience.

\subsection{MDEO-Cloud Solution}

\begin{enumerate}
    \item \textbf{Installation and Deployment (ID):} assessing setup effort, environment configuration, and ease of execution.
    \item \textbf{Documentation and Resources (DR):} evaluating documentation clarity, completeness, and usefulness of examples.
    \item \textbf{Metamodel and Transformation Implementation (MT):} measuring the effort required to model the problem, implement transformations, and express optimisation behaviour.
    \item \textbf{Objective and Constraint Implementation (OC):} assessing how easily objective functions and hard constraints can be expressed and maintained.
    \item \textbf{Diagnostics and Error Handling (DE):} rating error message quality, debugging support, and fault localisation.
    \item \textbf{Optimisation Configuration Usability (OU):} evaluating how easy it is to configure search parameters, solver settings, and experiment runs.
\end{enumerate}

\subsection{Python-Based Solution}

\textit{To be completed by the author of the Python-based solution.}

\reportsection{Limitations}

\subsection{General Limitations}

\begin{itemize}
    \item \textbf{Benchmarking scope:} the evaluation is limited to the IHTC 2024 instances. Their resource mix, planning horizons, and demand patterns may not represent every hospital environment.
    \item \textbf{Scope of conclusions:} the results compare two implementations on the selected IHTC benchmarks; they do not establish that either approach is superior for all healthcare scheduling settings.
    \item \textbf{Published-baseline comparability:} any comparison with published GIPS-based IHTC results is contextual only, because those results were obtained with a different solver paradigm, time-limit policy, and instance coverage.
    \item \textbf{Implementation differences:} the two solutions may use different internal representations, search operators, and objective decompositions. The reported measures must therefore be interpreted in their implementation-specific context.
    \item \textbf{Single-run sampling:} each instance is executed once because repeated runs of the larger instances are computationally expensive. The evaluation therefore reports observed outcomes rather than means, medians, confidence intervals, or statistical significance claims about stochastic-search reliability.
    \item \textbf{Model and conversion fidelity:} results depend on the correctness of the conversion from IHTC input data to each implementation and on the completeness of each model's constraints and objective calculations.
\end{itemize}

\subsection{MDEO-Cloud Solution}

\begin{itemize}
    \item \textbf{Runtime comparability:} MDEO-Cloud terminates after the fixed 300-generation limit rather than a wall-clock limit, regardless of whether the candidate set has converged. Runtime therefore measures the computational cost of one fixed-length search, not an identical amount of computation across implementations.
    \item \textbf{Test~10 execution conditions:} Test~10 used additional CPU threads. Its 60-second script and model-transformation safeguards were not reached and did not affect termination or the reported result. It remains useful as an annotated large-instance observation, but it is not a directly comparable runtime point for the Tests~01--09 scalability trend or for cross-framework runtime comparisons without equivalent execution conditions.
    \item \textbf{Fixed generation budget:} Test~10 reached the configured 300-generation limit without clear evidence of convergence. The reported Pareto candidates may therefore be inferior to solutions obtainable with a larger search budget.
\end{itemize}

\subsection{Python-Based Solution}

\textit{To be completed by the author of the Python-based solution.}

\begin{thebibliography}{9}

\bibitem{ceschia2025ihtc}
S.~Ceschia, R.~M.~Rosati, A.~Schaerf, P.~Smet, G.~Vanden Berghe, and E.~Zanazzo,
``The integrated healthcare timetabling competition 2024,''
\textit{Operations Research, Data Analytics and Logistics}, vol.~45, art.~200481, 2025.

\bibitem{deb2002nsga2}
K.~Deb, A.~Pratap, S.~Agarwal, and T.~Meyarivan,
``A fast and elitist multiobjective genetic algorithm: NSGA-II,''
\textit{IEEE Transactions on Evolutionary Computation}, vol.~6, no.~2, pp.~182--197, 2002.

\bibitem{krieger2026mdeo}
N.~Krieger,
\textit{A Cloud-Based Platform for Model-Driven Evolutionary Optimization},
Master's thesis, University of Stuttgart, 2026.

\bibitem{kratz2025gips}
M.~Kratz, S.~Zschaler, D.~Str\"uber, and T.~Goldschmidt,
\textit{High-Level Specification of MILP Problems with Class Diagrams and Graph Transformations},
SSRN preprint, 2025. Available at: \url{https://ssrn.com/abstract=5768119}.

\end{thebibliography}

\end{document}
